\section{Problem Formulation}

    Throughout our paper, a corresponding static-dynamic LiDAR scan pair, refers to a LiDAR scan pair $(s_i, d_i)$ where both the scans share exactly the same surrounding. The only difference is that $d_i$ consists of movable and moving objects, unlike $s_i$ (Figure \ref{fig:movable-first}). We have a set of static scans $\textbf{St}=(s_i, i: 0,1,2...)$ and corresponding dynamic scans denoted as $\textbf{Dy}=(d_i, i: 0,1,2...)$. Given a dynamic scan $d_i$, and its latent representation on the dynamic manifold, ($r_{d_i}$): our goal is to map $r_{d_i}$  to its closest corresponding point ($\overline{r}_{s_i}$) on the static manifold. $\overline{r}_{s_i}$ must produce the closest corresponding static LiDAR $\overline{s}_{i}$, by generating static background of the occluded regions, as seen in $s_i$. We achieve this by developing a novel discriminator and a contrastive strategy that allows precise discrimination of dynamic objects and occlusions from static objects and occlusions. The generated static structures along with the input dynamic LiDAR are used for movable and moving object segmentation.